\section{Design study}
\label{sec:results}

Every number in this section is produced by \code{paper/make\_results.py}
from the engine, together with the configuration that produced it
(\cref{sec:engine}). Unless stated otherwise the flights use a window of
\SI{\rTwindow}{\second} at a step of \SI{\rDt}{\milli\second}, statistics
taken after a settling interval of \SI{\rSettle}{\second}, gust seed
\rSeed, and the governor of \cref{mod:governor} at
$v_{\max}=\SI{2.5}{\metre\per\second}$,
$a_{\max}=\SI{2.5}{\metre\per\second\squared}$,
$\varepsilon=\SI{\rTrackMargin}{\milli\metre}$, heading rate cap
\SI{1.5}{\radian\per\second}.

\subsection{The requirement, worked backwards}

Four human-facing requirements bound the design.

\begin{enumerate}[label=(R\arabic*),nosep]
  \item \emph{Readability.} Two quantities, and they pull in opposite
        directions. Pixel density in pixels per degree,
        $1920/(2\arctan(w/2d))$ for a panel of width $w$ at distance $d$,
        measures sharpness and rises with distance. Text size is an angular
        height: ANSI/HFES 100~\cite{hfes2007} asks for a character height of
        at least \SI{16}{\arcminute}, with \SIrange{20}{22}{\arcminute}
        preferred, and it falls with distance. A desktop meets the second by
        scaling its interface, which divides the usable workspace by the
        scale factor.
  \item \emph{Noise.} A quiet office is \SI{45}{\decibel A}, ordinary speech
        \SI{60}{\decibel A}. Working beside a source for half an hour requires
        something in between.
  \item \emph{Downwash.} Loose paper begins to move at roughly
        \SI{5}{\metre\per\second}. By \cref{eq:dl} the wake velocity is
        $2\vind = 2\sqrt{\mathrm{DL}/2\rho}$, so this bounds disk loading.
  \item \emph{Size and safety.} The machine must fit beside a desk, and by
        \cref{eq:reach} its blades extend $\mathcal{R}$ from its centre.
\end{enumerate}

All four pull the same way, toward large slowly turning lightly loaded rotors,
and (R2) to (R4) conflict with the span limit through $\mathcal{R}=L+R$.

\subsection{The binding constraint is acoustic}

Applying \cref{eq:mmaxexplicit} to the fixed load of the design below
(\SI{\rFixed}{\gram} of panel, electronics and gimbal), the energetic wall is
comfortable at short sessions: at \SI{20}{\minute} the critical effective load
is \SI{\rWallTwenty}{\kilo\gram} against a requirement of
\SI{\rMzeroTwenty}{\kilo\gram}, a utilisation of \SI{\rWallUse}{\percent}.
Half an hour of flight costs about a quarter of a kilogram of cells. Energy
is not what stops this product.

Noise is. \Cref{obs:sixthpower} says the only lever with authority is tip
speed, and \cref{cor:designrule} says solidity is how one buys it. Optimising
under (R1) to (R4) drives the design to four rotors of
\SI{\rDiam}{\milli\metre} diameter with wide, low-aspect-ratio blades turning
at \SI{\rRpm}{rpm}, a tip speed of \SI{\rVtip}{\metre\per\second} against
roughly \SI{68}{\metre\per\second} for a commercial camera drone.

\subsection{Which display}

A resolution is not a payload. \Cref{tab:displays} evaluates the candidates at
the standoff the safety analysis below will impose, \SI{\rStandoff}{\metre},
each with the aircraft re-optimised around it by differential evolution over
rotor diameter, blade loading, aspect ratio and blade count under a span limit
of \SI{\rOptSpan}{\metre}, a free-field level of \SI{\rOptSpl}{\decibel A} at
\SI{1}{\metre}, a wake of \SI{\rOptWash}{\metre\per\second} and a gross mass
of \SI{\rOptMass}{\kilo\gram}. Every $1920$-wide panel is sharp at that
distance; none delivers a full desktop there, for the reason in (R1). Areal
density is what discriminates between the aircraft: a desktop monitor panel is
more than twice as heavy per unit area as a laptop panel, because it carries
thicker glass, a heavier backlight and a metal chassis.

\begin{table}[t]
\centering
\caption{Display options at \SI{\rStandoff}{\metre}, each with the aircraft
re-optimised under identical limits. ``Logical'' is the desktop width in
logical pixels once the interface is scaled to a \SI{16}{\arcminute} character
height. Areal density in \si{\kilo\gram\per\metre\squared}, n/a for a whole
device. Levels in \si{\decibel A}, free field at \SI{1}{\metre} and at the
user's ear in a $5\times4\times\SI{2.4}{\metre}$ room with $\bar\alpha=0.20$
by \cref{prop:room}. ``Over limits'' means the closure exists but the best
design found violates at least one of the four limits.}
\label{tab:displays}
\footnotesize\setlength{\tabcolsep}{4pt}
\begin{tabular}{@{}lrrrrrrrl@{}}
\toprule
Display & px/deg & Logical & Areal & Gross & Power & $L_p(\SI{1}{\metre})$ & At ear & Verdict\\
\midrule
\inputrows{results/tab_displays_rows}
\end{tabular}
\end{table}

\Cref{fig:displays} shows the same comparison graphically.

\begin{figure}[tbp]
\centering
\providecommand{\eDsafeCross}{1.06}
\providecommand{\eDtext}{0.43}
\providecommand{\eDfull}{0.65}
\providecommand{\eLogicalAtCross}{784}
\providecommand{\eArmL}{0.3827}
\providecommand{\eRotorR}{0.2460}
\providecommand{\eSpan}{1.257}
\providecommand{\eReach}{0.629}
\providecommand{\eScreenW}{0.345}
\providecommand{\eScreenH}{0.194}
\providecommand{\eRcp}{0.130}
\providecommand{\eEye}{1.500}
\providecommand{\eZrotor}{1.370}
\providecommand{\eZrotorBelow}{1.630}
\providecommand{\eStandoff}{1.20}
\providecommand{\eDsafe}{0.45}
\providecommand{\eDeltaZ}{0.130}
\providecommand{\eDisplayOrder}{13.3 in laptop panel,15.6 in laptop panel,iPad Pro 13 panel,iPad Pro 13 whole,15.6 in lid assembly,21.5 in monitor panel,24 in monitor panel,13.3 in ultrabook,15.6 in laptop whole}
\providecommand{\eOptSpl}{58}
\providecommand{\eOptMass}{6.0}
\providecommand{\eTurnClear}{0.61}
\providecommand{\eTurnTrack}{64}
\providecommand{\eTrackMargin}{100}
\providecommand{\eTurnPmean}{111.9}
\providecommand{\eHoverBat}{288}
\providecommand{\eReservePct}{12}
\providecommand{\eTethPsrc}{109.6}
\providecommand{\eTethArea}{0.228}
\providecommand{\eTethV}{48}
\providecommand{\eTethL}{5}
\providecommand{\eTethDrop}{5}
%
\begin{tikzpicture}
\begin{groupplot}[
  group style={group size=2 by 1, horizontal sep=0.35cm,
               yticklabels at=edge left},
  paperplot, width=0.40\linewidth, height=6.4cm,
  xbar, /pgf/bar width=6.5pt, /pgf/bar shift=0pt,
  y dir=reverse, symbolic y coords/.expand once={\eDisplayOrder},
  ytick/.expand once={\eDisplayOrder}, enlarge y limits=0.09,
  y tick label style={font=\footnotesize},
  legend style={at={(0.98,0.98)}, anchor=north east}]
\nextgroupplot[xlabel={gross mass $m$ [\si{\kilogram}]}, xmin=0, xmax=7]
  \addplot[fill=cA!70, draw=cA] table[x=gross, y=tag] {results/fig/e_displays_ok.dat};
  \addplot[fill=cD!55, draw=cD] table[x=gross, y=tag] {results/fig/e_displays_bad.dat};
  \draw[cD, dashed, line width=0.8pt]
    (axis cs:\eOptMass,{13.3 in laptop panel}) -- (axis cs:\eOptMass,{15.6 in laptop whole});
  \addlegendimage{line legend, cD, dashed, line width=0.8pt}
  \legend{closes, over limits, limit}
\nextgroupplot[xlabel={$L_p(\SI{1}{\metre})$ [dBA]}, xmin=45, xmax=70]
  \addplot[fill=cA!70, draw=cA] table[x=spl, y=tag] {results/fig/e_displays_ok.dat};
  \addplot[fill=cD!55, draw=cD] table[x=spl, y=tag] {results/fig/e_displays_bad.dat};
  \draw[cD, dashed, line width=0.8pt]
    (axis cs:\eOptSpl,{13.3 in laptop panel}) -- (axis cs:\eOptSpl,{15.6 in laptop whole});
\end{groupplot}
\end{tikzpicture}

\caption{The display candidates of \cref{tab:displays}, each with its aircraft
re-optimised. Left: gross mass, against the \SI{\eOptMass}{\kilo\gram}
limit. Right: free-field level at \SI{1}{\metre}, against the
\SI{\eOptSpl}{dBA} limit. Every candidate that fails exceeds the noise limit,
and only one of them exceeds the mass limit: a heavier display needs a more
heavily loaded rotor, and by \cref{obs:sixthpower} that is paid for in
decibels long before it is paid for in kilograms.}
\label{fig:displays}
\end{figure}

\begin{finding}
The full $1920\times1080$ panel is available, provided it is a laptop panel
and not a monitor panel. Flying a whole laptop is not: the panel is a sixth
of a 15.6-inch machine's mass and the rest is keyboard, battery, cooling and
chassis, none of which is visible while hovering. A tablet is different in
kind because a tablet is almost entirely screen, which is why the iPad flown
whole still closes while the laptop flown whole does not. What no panel of
this size delivers at \SI{\rStandoff}{\metre} is a full desktop's worth of
text: see \cref{sec:readability}.
\end{finding}

\subsection{Geometry relative to the user}

The standoff a user specifies is measured to the \emph{display}. The rotors
extend $\mathcal{R}=L+R$ back toward them from the machine's centre.

\begin{proposition}[Blade clearance]
\label{prop:clearance}
Let the display be at horizontal standoff $d$ from the user's eye, carried a
distance $b\ge0$ ahead of the rotor centroid on a boom, with the rotor plane a
height $\Delta z$ above the eye. Then the nearest blade tip lies at
\begin{equation}
  \ell_0 = \sqrt{\big(\max\{d+b-\mathcal{R},0\}\big)^2 + \Delta z^2} ,
  \label{eq:clearance}
\end{equation}
and the requirement $\ell_0\ge d_{\mathrm{safe}}$ is met either by
\begin{equation}
  d \ge \mathcal{R} + \sqrt{d_{\mathrm{safe}}^2 - \Delta z^2} - b
  \qquad\text{or}\qquad
  b \ge \mathcal{R} + \sqrt{d_{\mathrm{safe}}^2-\Delta z^2} - d .
  \label{eq:clearancefix}
\end{equation}
\end{proposition}

At a nominal $d=\SI{0.70}{\metre}$ with $\mathcal{R}=\SI{\rReach}{\metre}$
and $|\Delta z|=\SI{0.13}{\metre}$, \cref{eq:clearance} gives
$\ell_0=\SI{0.148}{\metre}$ against a safe distance of \SI{0.45}{\metre}: a
violation by a factor of three, invisible to any check evaluated at the
vehicle centre.

\Cref{fig:clearance} draws the two standoffs to scale.

\begin{figure}[tbp]
\centering
\providecommand{\eDsafeCross}{1.06}
\providecommand{\eDtext}{0.43}
\providecommand{\eDfull}{0.65}
\providecommand{\eLogicalAtCross}{784}
\providecommand{\eArmL}{0.3827}
\providecommand{\eRotorR}{0.2460}
\providecommand{\eSpan}{1.257}
\providecommand{\eReach}{0.629}
\providecommand{\eScreenW}{0.345}
\providecommand{\eScreenH}{0.194}
\providecommand{\eRcp}{0.130}
\providecommand{\eEye}{1.500}
\providecommand{\eZrotor}{1.370}
\providecommand{\eZrotorBelow}{1.630}
\providecommand{\eStandoff}{1.20}
\providecommand{\eDsafe}{0.45}
\providecommand{\eDeltaZ}{0.130}
\providecommand{\eDisplayOrder}{13.3 in laptop panel,15.6 in laptop panel,iPad Pro 13 panel,iPad Pro 13 whole,15.6 in lid assembly,21.5 in monitor panel,24 in monitor panel,13.3 in ultrabook,15.6 in laptop whole}
\providecommand{\eOptSpl}{58}
\providecommand{\eOptMass}{6.0}
\providecommand{\eTurnClear}{0.61}
\providecommand{\eTurnTrack}{64}
\providecommand{\eTrackMargin}{100}
\providecommand{\eTurnPmean}{111.9}
\providecommand{\eHoverBat}{288}
\providecommand{\eReservePct}{12}
\providecommand{\eTethPsrc}{109.6}
\providecommand{\eTethArea}{0.228}
\providecommand{\eTethV}{48}
\providecommand{\eTethL}{5}
\providecommand{\eTethDrop}{5}
%
\begin{tikzpicture}[x=3.3cm, y=3.3cm]
\foreach \dd/\sx/\col/\ellv/\ttl in {
    0.70/0.00/cD/0.148/{$d=\SI{0.70}{\metre}$, no boom: violated},
    1.20/2.05/cC/0.586/{$d=\SI{1.20}{\metre}$, no boom: the design}} {
  \begin{scope}[shift={(\sx,0)}]
    \pgfmathsetmacro\tip{\dd-\eReach}
    \pgfmathsetmacro\zlo{\eEye-\eScreenH/2}
    \pgfmathsetmacro\zhi{\eEye+\eScreenH/2}
    % the safe sphere about the eye
    \fill[cD!6] (0,\eEye) circle (\eDsafe);
    \draw[cD, dashed, line width=0.6pt] (0,\eEye) circle (\eDsafe);
    \node[lbl, text=cD, anchor=north] at (0,\eEye-\eDsafe-0.01) {$d_{\mathrm{safe}}$ about the eye};
    % the user
    \draw[line width=0.8pt, fill=white] (-0.10,\eEye+0.02) circle (0.10);
    \draw[line width=0.8pt] (-0.10,\eEye-0.08) -- (-0.10,1.12);
    \fill (0,\eEye) circle (0.9pt);
    \node[lbl, anchor=south] at (-0.10,\eEye+0.12) {eye};
    % rotor plane seen from the side: arms, two disks edge on, hubs
    \draw[line width=1.0pt] (\dd-\eArmL,\eZrotor) -- (\dd+\eArmL,\eZrotor);
    \draw[cB, line width=2.6pt] (\dd-\eArmL-\eRotorR,\eZrotor) -- (\dd-\eArmL+\eRotorR,\eZrotor);
    \draw[cB, line width=2.6pt] (\dd+\eArmL-\eRotorR,\eZrotor) -- (\dd+\eArmL+\eRotorR,\eZrotor);
    \fill (\dd-\eArmL,\eZrotor) circle (0.7pt) (\dd+\eArmL,\eZrotor) circle (0.7pt);
    % mast and display, screen above the rotors
    \draw[cG, line width=1.0pt] (\dd,\eZrotor) -- (\dd,\zlo);
    \draw[cA, line width=2.4pt] (\dd,\zlo) -- (\dd,\zhi);
    \node[lbl, text=cA, anchor=west] at (\dd+0.03,\zhi-0.02) {display};
    % the nearest tip and the clearance
    \draw[\col, line width=1.1pt] (0,\eEye) -- (\tip,\eZrotor);
    \fill[\col] (\tip,\eZrotor) circle (1.2pt);
    \node[lbl, text=\col!80!black, anchor=south west, inner sep=1pt]
      at (\tip+0.04,\eZrotor+0.035) {$\ell_0=\SI{\ellv}{\metre}$};
    % dimensions
    \draw[dim] (0,1.86) -- (\dd,1.86) node[midway, above, lbl] {$d$};
    \draw[construct] (\dd,\zhi) -- (\dd,1.89);
    \draw[construct] (0,\eEye) -- (0,1.89);
    \draw[dim] (\tip,1.24) -- (\dd,1.24) node[midway, below, lbl] {$\mathcal{R}=L+R$};
    \draw[construct] (\tip,\eZrotor) -- (\tip,1.21);
    \draw[construct] (\dd,\eZrotor) -- (\dd,1.21);
    \draw[dim] (-0.33,\eZrotor) -- (-0.33,\eEye) node[midway, left, lbl] {$\Delta z$};
    \draw[construct] (-0.36,\eZrotor) -- (\tip,\eZrotor);
    \draw[construct] (-0.36,\eEye) -- (-0.21,\eEye);
    \node[font=\small, anchor=south] at (\dd/2+0.25,2.00) {\ttl};
  \end{scope}
}
\end{tikzpicture}

\caption{\Cref{prop:clearance} in side view, to scale. The display is held at
eye height at standoff $d$; with the screen above the rotors the rotor plane
lies $|\Delta z|=\SI{\eDeltaZ}{\metre}$ below the eye; the nearest blade tip is
the reach $\mathcal{R}=L+R$ back toward the user from the rotor centroid. At
$d=\SI{0.70}{\metre}$ the tip is inside the sphere of radius $d_{\mathrm{safe}}$
about the eye, although the display is well outside it; a check made at the
vehicle centre would pass this design. At $d=\SI{1.20}{\metre}$ the tip is
outside.}
\label{fig:clearance}
\end{figure}

\Cref{tab:clearance} compares the two remedies, each option closed, its gains
capped by \cref{mod:gains} for its own authority (attitude and position gains
together, since the cascade needs its time-scale separation), and flown
through the turnaround mission with the governor active. For every option but
the last standoff the ideal display position lies inside the keep-out radius
$\varrho = d_{\mathrm{safe}}+\mathcal{R}+\varepsilon$, by up to
\SI{\rCldNearInside}{\centi\metre}; those missions start on the sphere and the
governor holds the reference there, so their ``turn'' clearance measures the
keep-out sphere and not the option. The sphere is conservative for a boom,
whose reach is set by the rear rotors while the front rotors are what
approach the user.

\begin{table}[t]
\centering
\caption{Two ways to recover blade clearance. The boom keeps the display close
and pays in pitch inertia and authority; standing back is free on every axis
and quieter, and pays in text size. ``Turn'' is the minimum distance from the
eye to the nearest rotor disk during the turnaround mission, against the
\SI{0.45}{\metre} required. $J_{yy}$ in \si{\kilo\gram\metre\squared},
$\alpha^{\max}_{\mathrm{pitch}}$ in \si{\radian\per\second\squared}, $K_R$ in
\si{\per\second\squared}, screen angle is the panel's apparent width, logical
width as in \cref{tab:displays}.}
\label{tab:clearance}
\footnotesize\setlength{\tabcolsep}{4pt}
\begin{tabular}{@{}lrrrrrrrrr@{}}
\toprule
Option & Tip gap & Gross & Power & $J_{yy}$ & $\alpha^{\max}_{\mathrm{pitch}}$
       & $K_R$ & Screen & Logical & Turn\\
\midrule
\inputrows{results/tab_clearance_rows}
\end{tabular}
\end{table}

\begin{finding}
A boom of \SI{0.38}{\metre} multiplies the pitch inertia by \rBoomJratio,
reduces the pitch authority to \rBoomAlphaRatio\ of the no-boom value and the
admissible attitude gain to $K_R=\rClbMidKR$, and costs \SI{\rBoomGross}{\gram}
of gross mass and \SI{\rBoomEar}{\decibel A} at the ear against
\SI{\rStandEar}{\decibel A}. Standing back to \SI{1.20}{\metre} achieves the
same static clearance at no cost in mass, power or authority, is the only
option whose ideal position is outside the keep-out sphere, and keeps its
rotors \SI{\rCldFarClear}{\metre} from the eye through the turnaround. The
price is the one (R1) names, and it is paid below.
\end{finding}

\subsection{What the user can read}
\label{sec:readability}

At \SI{\rStandoff}{\metre} the 15.6-inch panel subtends \ang{\rScreenDeg} at
\SI{\rPpd}{px\per\degree}: sharper than any desktop monitor. But a
\SI{16}{px} body-text character, cap height \SI{11}{px}, subtends only
\SI{\rCapArcmin}{\arcminute} at that distance, against the \SI{16}{\arcminute}
minimum of~\cite{hfes2007}. Meeting the minimum takes an interface scale of
\SI{\rUiScale}{\percent} (\SI{\rUiScalePref}{\percent} for the preferred
\SI{20}{\arcminute}), which leaves a logical desktop of
$\rLogicalW\times\rLogicalH$ pixels. That is a tablet's workspace on a
laptop's panel. At \SI{0.70}{\metre} the same panel gives
\SI{\rCapArcminNear}{\arcminute} and a logical width of \rLogicalWNear, a
laptop's workspace, and its blades \SI{0.15}{\metre} from the eye.

\begin{finding}
The two safe ways to fly a panel this size, standing back or a boom, and the
requirement that the user read a full desktop on it, are not simultaneously
satisfiable at $1920\times1080$ in \SI{345}{\milli\metre}. A safe flying
screen at reading distance either carries a physically larger panel, which
\cref{tab:displays} shows the acoustics will not allow, or presents a smaller
workspace than the resolution suggests. The resolution is not the constraint;
the angular size of a character is.
\end{finding}

\Cref{fig:readability} shows the whole trade over standoff.

\begin{figure}[tbp]
\centering
\providecommand{\eDsafeCross}{1.06}
\providecommand{\eDtext}{0.43}
\providecommand{\eDfull}{0.65}
\providecommand{\eLogicalAtCross}{784}
\providecommand{\eArmL}{0.3827}
\providecommand{\eRotorR}{0.2460}
\providecommand{\eSpan}{1.257}
\providecommand{\eReach}{0.629}
\providecommand{\eScreenW}{0.345}
\providecommand{\eScreenH}{0.194}
\providecommand{\eRcp}{0.130}
\providecommand{\eEye}{1.500}
\providecommand{\eZrotor}{1.370}
\providecommand{\eZrotorBelow}{1.630}
\providecommand{\eStandoff}{1.20}
\providecommand{\eDsafe}{0.45}
\providecommand{\eDeltaZ}{0.130}
\providecommand{\eDisplayOrder}{13.3 in laptop panel,15.6 in laptop panel,iPad Pro 13 panel,iPad Pro 13 whole,15.6 in lid assembly,21.5 in monitor panel,24 in monitor panel,13.3 in ultrabook,15.6 in laptop whole}
\providecommand{\eOptSpl}{58}
\providecommand{\eOptMass}{6.0}
\providecommand{\eTurnClear}{0.61}
\providecommand{\eTurnTrack}{64}
\providecommand{\eTrackMargin}{100}
\providecommand{\eTurnPmean}{111.9}
\providecommand{\eHoverBat}{288}
\providecommand{\eReservePct}{12}
\providecommand{\eTethPsrc}{109.6}
\providecommand{\eTethArea}{0.228}
\providecommand{\eTethV}{48}
\providecommand{\eTethL}{5}
\providecommand{\eTethDrop}{5}
%
\begin{tikzpicture}
\begin{groupplot}[
  group style={group size=1 by 2, vertical sep=0.35cm,
               xlabels at=edge bottom, xticklabels at=edge bottom},
  paperplot, height=4.6cm, xmin=0.35, xmax=2.0,
  xlabel={standoff $d$ from the eye to the display [\si{\metre}]}]
% (a) usable desktop once text is scaled to 16 arcmin
\nextgroupplot[ylabel={logical width [px]}, ymin=0, ymax=2150,
  ytick={0,500,1000,1280,1920}]
  \fill[cC!14] (axis cs:0.35,0) rectangle (axis cs:\eDfull,2150);
  \addplot[cA] table[x=d, y=logical] {results/fig/e_readability.dat};
  \addplot[cG, dashed, domain=0.35:2.0, samples=2] {1280};
  \node[lbl, cG, anchor=south east] at (axis cs:1.98,1280) {full desktop, \num{1280} px};
  \node[lbl, text=cC!55!black, anchor=south west, align=left]
    at (axis cs:0.37,120) {full\\ desktop\\ legible};
  \draw[construct] (axis cs:\eDsafeCross,0) -- (axis cs:\eDsafeCross,2150);
  \draw[cA, densely dotted, line width=0.7pt] (axis cs:\eStandoff,0) -- (axis cs:\eStandoff,2150);
  \node[lbl, text=cA, anchor=south west] at (axis cs:\eStandoff,1500) {design};
  \addplot[only marks, mark=*, mark size=1.8pt, cB]
    coordinates {(\eDsafeCross,\eLogicalAtCross)};
  \node[lbl, text=cB!80!black, anchor=south west]
    at (axis cs:\eDsafeCross,\eLogicalAtCross) {\eLogicalAtCross\ px at $d=\SI{\eDsafeCross}{\metre}$};
% (b) nearest blade tip, prop:clearance
\nextgroupplot[ylabel={blade tip to eye $\ell_0$ [\si{\metre}]}, ymin=0, ymax=1.35]
  \fill[cC!14] (axis cs:\eDsafeCross,0) rectangle (axis cs:2.0,1.35);
  \addplot[cB] table[x=d, y=gap] {results/fig/e_readability.dat};
  \addplot[cD, dashed, domain=0.35:2.0, samples=2] {\eDsafe};
  \node[lbl, text=cD, anchor=south west] at (axis cs:0.37,\eDsafe) {$d_{\mathrm{safe}}=\SI{\eDsafe}{\metre}$};
  \node[lbl, text=cC!55!black, anchor=north east, align=right]
    at (axis cs:1.98,1.30) {blades outside\\ $d_{\mathrm{safe}}$};
  \node[lbl, text=cB!80!black, anchor=south west, align=left]
    at (axis cs:0.37,0.16) {rotors overhang the user: $\ell_0=|\Delta z|$};
  \draw[construct] (axis cs:\eDsafeCross,0) -- (axis cs:\eDsafeCross,1.35);
  \draw[cA, densely dotted, line width=0.7pt] (axis cs:\eStandoff,0) -- (axis cs:\eStandoff,1.35);
\end{groupplot}
\end{tikzpicture}

\caption{Readability against blade clearance, over the standoff $d$ from the
eye to the display. Top: the logical desktop width once the interface is
scaled to a \SI{16}{\arcminute} character height; it is the full \num{1920}
pixels only within \SI{\eDtext}{\metre}, where unscaled text already meets the
minimum, and it falls as $1/d$ beyond. A \num{1280}-pixel desktop needs
$d\le\SI{\eDfull}{\metre}$ (green). Bottom: the nearest blade tip $\ell_0$ of
\cref{eq:clearance}; it stays at $|\Delta z|$ while the rotors overhang the
user and clears $d_{\mathrm{safe}}$ only for $d\ge\SI{\eDsafeCross}{\metre}$
(green). The two green regions do not overlap, and at the first safe standoff
the desktop is already down to \eLogicalAtCross\ pixels.}
\label{fig:readability}
\end{figure}

\subsection{Screen above or below the rotors}

Because the display is vertical and the rotor flow is vertical
(\cref{sec:dynamics}), placing the panel above or below the rotor plane is
nearly free aerodynamically: \SI{\rGrossBelow}{\kilo\gram} against
\SI{\rGross}{\kilo\gram}, with the same noise. What changes is geometry. With
the screen below, the rotor plane sits at \SI{\rZrotorBelow}{\metre}, level
with the user's head; with it above, the rotors drop to
\SI{\rZrotorAbove}{\metre} and the panel itself lies between the blades and
the user's face. The latter is adopted. Tracking is marginally better in still
air, because by the remark following \cref{eq:screenmoment} the drag moment
becomes weathercock-stable, and marginally worse in a steady breeze for the
same reason.

\Cref{fig:placement} compares the two at the design standoff.

\begin{figure}[tbp]
\centering
\providecommand{\eDsafeCross}{1.06}
\providecommand{\eDtext}{0.43}
\providecommand{\eDfull}{0.65}
\providecommand{\eLogicalAtCross}{784}
\providecommand{\eArmL}{0.3827}
\providecommand{\eRotorR}{0.2460}
\providecommand{\eSpan}{1.257}
\providecommand{\eReach}{0.629}
\providecommand{\eScreenW}{0.345}
\providecommand{\eScreenH}{0.194}
\providecommand{\eRcp}{0.130}
\providecommand{\eEye}{1.500}
\providecommand{\eZrotor}{1.370}
\providecommand{\eZrotorBelow}{1.630}
\providecommand{\eStandoff}{1.20}
\providecommand{\eDsafe}{0.45}
\providecommand{\eDeltaZ}{0.130}
\providecommand{\eDisplayOrder}{13.3 in laptop panel,15.6 in laptop panel,iPad Pro 13 panel,iPad Pro 13 whole,15.6 in lid assembly,21.5 in monitor panel,24 in monitor panel,13.3 in ultrabook,15.6 in laptop whole}
\providecommand{\eOptSpl}{58}
\providecommand{\eOptMass}{6.0}
\providecommand{\eTurnClear}{0.61}
\providecommand{\eTurnTrack}{64}
\providecommand{\eTrackMargin}{100}
\providecommand{\eTurnPmean}{111.9}
\providecommand{\eHoverBat}{288}
\providecommand{\eReservePct}{12}
\providecommand{\eTethPsrc}{109.6}
\providecommand{\eTethArea}{0.228}
\providecommand{\eTethV}{48}
\providecommand{\eTethL}{5}
\providecommand{\eTethDrop}{5}
%
% Side views at the design standoff, to scale: the display is held at eye
% height either way; what moves is the rotor plane.
\begin{tikzpicture}[x=3.0cm, y=3.0cm]
\foreach \zr/\sx/\anc/\dy/\ttl in {
    \eZrotorBelow/0.00/south east/0.03/{screen below the rotors},
    \eZrotor/2.30/north east/-0.03/{screen above the rotors: adopted}} {
  \begin{scope}[shift={(\sx,0)}]
    \pgfmathsetmacro\tip{\eStandoff-\eReach}
    \pgfmathsetmacro\zlo{\eEye-\eScreenH/2}
    \pgfmathsetmacro\zhi{\eEye+\eScreenH/2}
    % eye height across the panel
    \draw[construct] (-0.22,\eEye) -- (\eStandoff+\eReach,\eEye);
    \node[lbl, text=cG!80!black, anchor=south west] at (-0.22,\eEye+0.13) {eye height \SI{\eEye}{\metre}};
    \draw[construct, -{Stealth[length=1.2mm]}] (-0.02,\eEye+0.14) -- (-0.02,\eEye+0.01);
    % the user: head and shoulders
    \draw[line width=0.8pt, fill=white] (-0.10,\eEye+0.02) circle (0.10);
    \draw[line width=0.8pt] (-0.10,\eEye-0.08) -- (-0.10,1.10);
    \draw[line width=0.8pt] (-0.24,1.36) -- (0.04,1.36);
    \fill (0,\eEye) circle (0.9pt);
    % rotor plane: arms, two disks edge on, hubs
    \draw[line width=1.0pt] (\eStandoff-\eArmL,\zr) -- (\eStandoff+\eArmL,\zr);
    \draw[cB, line width=2.6pt] (\eStandoff-\eArmL-\eRotorR,\zr) -- (\eStandoff-\eArmL+\eRotorR,\zr);
    \draw[cB, line width=2.6pt] (\eStandoff+\eArmL-\eRotorR,\zr) -- (\eStandoff+\eArmL+\eRotorR,\zr);
    \fill (\eStandoff-\eArmL,\zr) circle (0.7pt) (\eStandoff+\eArmL,\zr) circle (0.7pt);
    \node[lbl, text=cB!80!black, anchor=\anc] at (\eStandoff+\eReach,\zr+\dy)
      {rotor plane \SI{\zr}{\metre}};
    % mast and display at eye height
    \draw[cG, line width=1.0pt] (\eStandoff,\zr) -- (\eStandoff,\eEye);
    \draw[cA, line width=2.4pt] (\eStandoff,\zlo) -- (\eStandoff,\zhi);
    \node[lbl, text=cA, anchor=west] at (\eStandoff+0.03,\zhi-0.03) {display};
    % the nearest tip
    \draw[cC, line width=1.1pt] (0,\eEye) -- (\tip,\zr);
    \fill[cC] (\tip,\zr) circle (1.2pt);
    \node[font=\small, anchor=south] at (\eStandoff/2+0.25,1.86) {\ttl};
  \end{scope}
}
\end{tikzpicture}

\caption{The two placements of the display at the design standoff, to scale,
with the display held at eye height in both. The nearest-tip clearance
(green) is the same either way, because \cref{eq:clearance} depends on
$\Delta z$ only through $\Delta z^2$. What differs is where the rotor plane
sits relative to the head: with the screen below it is at the top of the head,
with the screen above it drops below the chin, and the panel is then the
object between the blades and the face. The heights marked are rotor plane
heights above the floor.}
\label{fig:placement}
\end{figure}

\subsection{The design}

\Cref{tab:design} gives the resulting machine with a buildable payload: a bare
laptop panel on a matched controller board, a Raspberry Pi running a remote
desktop client, a separate flight controller, and an on-sensor inference
camera. \Cref{fig:design} draws it to scale.

\begin{figure}[tbp]
\centering
\providecommand{\eDsafeCross}{1.06}
\providecommand{\eDtext}{0.43}
\providecommand{\eDfull}{0.65}
\providecommand{\eLogicalAtCross}{784}
\providecommand{\eArmL}{0.3827}
\providecommand{\eRotorR}{0.2460}
\providecommand{\eSpan}{1.257}
\providecommand{\eReach}{0.629}
\providecommand{\eScreenW}{0.345}
\providecommand{\eScreenH}{0.194}
\providecommand{\eRcp}{0.130}
\providecommand{\eEye}{1.500}
\providecommand{\eZrotor}{1.370}
\providecommand{\eZrotorBelow}{1.630}
\providecommand{\eStandoff}{1.20}
\providecommand{\eDsafe}{0.45}
\providecommand{\eDeltaZ}{0.130}
\providecommand{\eDisplayOrder}{13.3 in laptop panel,15.6 in laptop panel,iPad Pro 13 panel,iPad Pro 13 whole,15.6 in lid assembly,21.5 in monitor panel,24 in monitor panel,13.3 in ultrabook,15.6 in laptop whole}
\providecommand{\eOptSpl}{58}
\providecommand{\eOptMass}{6.0}
\providecommand{\eTurnClear}{0.61}
\providecommand{\eTurnTrack}{64}
\providecommand{\eTrackMargin}{100}
\providecommand{\eTurnPmean}{111.9}
\providecommand{\eHoverBat}{288}
\providecommand{\eReservePct}{12}
\providecommand{\eTethPsrc}{109.6}
\providecommand{\eTethArea}{0.228}
\providecommand{\eTethV}{48}
\providecommand{\eTethL}{5}
\providecommand{\eTethDrop}{5}
%
\begin{tikzpicture}[x=4.9cm, y=4.9cm]
\pgfmathsetmacro\cc{\eArmL*cos(45)}
\pgfmathtruncatemacro\hubmm{round(2*\cc*1000)}
% ---- top view --------------------------------------------------------------
\begin{scope}
  \foreach \sx/\sy/\ang in {1/1/20, -1/1/-35, -1/-1/60, 1/-1/-10} {
    \coordinate (h) at (\sx*\cc,\sy*\cc);
    \draw[line width=1.6pt, cG!70!black] (0,0) -- (h);
    \draw[cB, fill=cB!7, line width=0.6pt] (h) circle (\eRotorR);
    \draw[cG, line width=0.4pt] (h) circle (\eRotorR+0.012);
    \draw[cB!70!black, line width=1.3pt]
      ($(h)+(\ang:\eRotorR-0.01)$) -- ($(h)+(\ang+180:\eRotorR-0.01)$);
    \fill (h) circle (0.8pt);
  }
  \draw[fill=cG!25] (-0.055,-0.055) rectangle (0.055,0.055);
  % the display seen from above: the panel lies in the b2-b3 plane
  \draw[cA, line width=2.4pt] (0,-\eScreenW/2) -- (0,\eScreenW/2);
  \draw[cA, line width=0.4pt] (0,-\eScreenW/2) -- (0,-0.62)
    node[below right, lbl, text=cA, xshift=-2pt] {display, above the rotor plane};
  % hub spacing across the top
  \draw[construct] (-\cc,\cc) -- (-\cc,\cc+\eRotorR+0.07);
  \draw[construct] (\cc,\cc) -- (\cc,\cc+\eRotorR+0.07);
  \draw[dim] (-\cc,\cc+\eRotorR+0.05) -- (\cc,\cc+\eRotorR+0.05)
    node[midway, above, lbl] {$2L\cos\ang{45}=\SI{\hubmm}{\milli\metre}$};
  % rotor diameter on rotor 1
  \draw[construct] (\cc,\cc+\eRotorR) -- (\cc+\eRotorR+0.08,\cc+\eRotorR);
  \draw[construct] (\cc,\cc-\eRotorR) -- (\cc+\eRotorR+0.08,\cc-\eRotorR);
  \draw[dim] (\cc+\eRotorR+0.06,\cc-\eRotorR) -- (\cc+\eRotorR+0.06,\cc+\eRotorR)
    node[midway, right, lbl] {$D=\SI{\rDiam}{\milli\metre}$};
  % body axes, set apart: b1 toward the user
  \draw[vec, cA] (-0.68,-0.64) -- (-0.50,-0.64) node[right, lbl] {$\vect b_1$};
  \draw[vec, cA] (-0.68,-0.64) -- (-0.68,-0.46) node[above, lbl] {$\vect b_2$};
  \node[lbl, anchor=west] at (0.60,-0.10) {user $\rightarrow$};
  \node[font=\small, anchor=south] at (0,0.70) {top view};
\end{scope}
% ---- side view --------------------------------------------------------------
\begin{scope}[shift={(1.50,0)}]
  \pgfmathsetmacro\zlo{\eRcp-\eScreenH/2}
  \pgfmathsetmacro\zhi{\eRcp+\eScreenH/2}
  \draw[line width=1.6pt, cG!70!black] (-\cc,0) -- (\cc,0);
  \draw[cB, line width=2.6pt] (-\cc-\eRotorR,0) -- (-\cc+\eRotorR,0);
  \draw[cB, line width=2.6pt] (\cc-\eRotorR,0) -- (\cc+\eRotorR,0);
  \draw[fill=cG!25] (-0.055,-0.035) rectangle (0.055,0.02);
  \draw[cG, line width=1.0pt] (0,0.02) -- (0,\zlo);
  \draw[cA, line width=2.4pt] (0,\zlo) -- (0,\zhi);
  \draw[cA, fill=white] (0,\eRcp) circle (0.011);
  \node[lbl, text=cA, anchor=west] at (0.02,\zhi) {display, gimbal pivot};
  \draw[vec, cA] (0.40,-0.20) -- (0.58,-0.20) node[right, lbl] {$\vect b_1$};
  \draw[vec, cA] (0.40,-0.20) -- (0.40,-0.02) node[above, lbl] {$\vect b_3$};
  % dimensions
  \draw[dim] (-0.16,0) -- (-0.16,\eRcp) node[midway, left, lbl] {$r_{cp}$};
  \draw[construct] (-0.19,\eRcp) -- (-0.012,\eRcp);
  \draw[dim] (0.12,\zlo) -- (0.12,\zhi) node[midway, right, lbl] {$h=\SI{194}{\milli\metre}$};
  \draw[dim] (-\cc-\eRotorR,-0.10) -- (\cc+\eRotorR,-0.10)
    node[midway, below, lbl] {$2(L\cos\ang{45}+R)$};
  \draw[construct] (-\cc-\eRotorR,0) -- ++(0,-0.12);
  \draw[construct] (\cc+\eRotorR,0) -- ++(0,-0.12);
  \node[font=\small, anchor=south] at (0,0.70) {side view};
\end{scope}
\end{tikzpicture}

\caption{DESK-1/W to scale. Top view: four rotors of diameter $D$ in the
$\times$ layout, the arm length $L$ set by the tip gap rule of
\cref{sec:components} so that adjacent guards (outer circles) nearly touch,
the display edge on above the hub, and $\vect b_1$ pointing at the user. Side
view: the display on its gimbal a height $r_{cp}$ above the rotor plane,
which is what puts the panel, and not the blades, between the rotors and the
user's face. The largest dimension, tip to tip across the diagonal, is the
span $2(L+R)=\SI{\rSpan}{\metre}$.}
\label{fig:design}
\end{figure}

\begin{table}[t]
\centering
\caption{DESK-1/W with the thin-client payload, screen above the rotors,
standoff \SI{\rStandoff}{\metre}. Sized for hover by \cref{eq:fpclosure}.}
\label{tab:design}
\footnotesize\setlength{\tabcolsep}{4pt}
\begin{tabular}{@{}llrl@{}}
\toprule
& Quantity & Value & Note\\
\midrule
\multirow{4}{*}{Geometry}
& rotors & $4\times\SI{\rDiam}{\milli\metre}$ & 2 blades, \SI{\rChord}{\milli\metre} chord, AR 3.0\\
& solidity & \rSolidity & below the 0.25 cascade limit\\
& span & \SI{\rSpan}{\metre} & reach \SI{\rReach}{\metre}\\
& display & $345\times\SI{194}{\milli\metre}$ & 15.6 in, \SI{\rPpd}{px\per\degree}, $\rLogicalW\times\rLogicalH$ logical\\
\midrule
\multirow{5}{*}{Mass}
& fixed payload & \SI{\rFixed}{\gram} & panel \SI{\rPanel}{\gram}, electronics \SI{\rElectronics}{\gram}, gimbal \SI{\rGimbal}{\gram}\\
& frame & \SI{\rFrame}{\gram} & arms \SI{\rArms}{\gram}, guards \SI{\rGuards}{\gram}\\
& propulsion & \SI{\rProp}{\gram} & motors \SI{\rMotors}{\gram}, props \SI{\rProps}{\gram}\\
& battery & \SI{\rBattery}{\gram} & \rPackS S1P \rPackCell, \SI{\rPackWh}{\watt\hour}\\
& \textbf{gross} & \textbf{\SI{\rGross}{\kilo\gram}} & $\FM=\rFM$ implied\\
\midrule
\multirow{4}{*}{Rotor}
& speed & \SI{\rRpm}{rpm} & \SI{\rRpmMax}{rpm} at $T_{\max}$\\
& tip speed & \SI{\rVtip}{\metre\per\second} & \SI{\rVtipMax}{\metre\per\second} at $T_{\max}$\\
& blade loading & \rCtSigma & stall at 0.14\\
& Reynolds & \rReynolds & $\Cdz = \rCdzero$ corrected\\
\midrule
\multirow{3}{*}{Power}
& hover & \SI{\rPower}{\watt} & of which \SI{\rPaux}{\watt} auxiliary\\
& session, hovering & \SI{\rEndReserve}{\minute} & to the \SI{12}{\percent} reserve; \SI{\rEndEmpty}{\minute} to empty\\
& session, worst mission & \SI{\rEndTurn}{\minute} & turnaround, to the reserve\\
\midrule
\multirow{4}{*}{Human factors}
& $L_p$ free field \SI{1}{\metre} & \SI{\rSplOne}{\decibel A} & \\
& $L_p$ at the ear, in room & \SI{\rSplEar}{\decibel A} & \cref{prop:room}, $r_c=\SI{\rCritDist}{\metre}$\\
& $L_p$ elsewhere in room & \SI{\rSplFar}{\decibel A} & does not fall with distance\\
& wake velocity & \SI{\rWake}{\metre\per\second} & Beaufort 3\\
\midrule
\multirow{4}{*}{Control}
& $J_{xx},J_{yy},J_{zz}$ & $\rJxx,\ \rJyy,\ \rJzz$ & \si{\kilo\gram\metre\squared}, about the centre of mass\\
& authority roll/pitch/yaw & $\rAlphaRoll/\rAlphaPitch/\rAlphaYaw$ & \si{\radian\per\second\squared}, \cref{eq:authority}\\
& gains & $K_R=\rKR$, $K_\omega=\rKw$ & below the bound \rKRsug\ of \cref{mod:gains}\\
& blade tip to face, static & \SI{\rTipGap}{\metre} & against \SI{0.45}{\metre} required\\
\bottomrule
\end{tabular}
\end{table}

\subsection{Flight performance}

\Cref{tab:flight} reports the design of \cref{tab:design}, with its hover-sized
battery, flown through each mission with the governor of \cref{mod:governor}
active. Two clearances are reported. The envelope gap is the distance from the
eye to a sphere of radius $\mathcal{R}$ about the centre of mass, the
conservative quantity that \cref{prop:governor}(iv) bounds. The rotor
clearance is the distance from the eye to the nearest oriented rotor disk,
which is what a blade could actually touch.

The six motions are not defined elsewhere, so they are stated here. Each is a
smooth trajectory of the person's position and heading over the
\SI{\rTwindow}{\second} window, and the ideal vehicle position follows from it
by holding the display at the standoff in front of the eyes. Standing is
postural sway of a few centimetres; pacing is walking back and forth along one
line while facing the display; sit to stand lowers the eyes by
\SI{0.45}{\metre}; walking a loop and jogging follow a circle of radius
\SI{1.44}{\metre}, jogging adding a vertical bob at the stride frequency; and
turning around reverses the heading at each end of a walk.
\Cref{fig:missions} draws them, together with the acceleration each one asks
of the display.

\begin{figure}[tbp]
\centering
% The six motions of the person, and where they put the display. Grey is the
% person, blue is the ideal vehicle position that holds the display at the
% standoff in front of their eyes, dashed arrows join the two.
\begin{tikzpicture}
\begin{groupplot}[
  group style={group size=3 by 2, horizontal sep=1.55cm, vertical sep=1.55cm},
  paperplot, width=0.325\linewidth, height=0.31\linewidth,
  title style={font=\small, yshift=-2pt},
  every axis plot/.append style={line width=0.9pt},
  xlabel style={font=\footnotesize}, ylabel style={font=\footnotesize},
  legend style={font=\scriptsize},
  topview/.style={axis equal image, xmin=-2.3, xmax=2.3, ymin=-2.3, ymax=2.3,
    xlabel={$x$ [\si{\metre}]}, ylabel={$y$ [\si{\metre}]}},
  timeaxis/.style={xmin=0, xmax=16, xlabel={$t$ [\si{\second}]}},
  person/.style={cG!60, line width=2.4pt},
  standoff/.style={construct, -{Stealth[length=1.3mm]}, each nth point=40,
    quiver={u=\thisrow{dx}, v=\thisrow{dy}}},
  stations/.style={only marks, mark=*, mark size=1.1pt, cG!80!black, each nth point=40}]
% row 1: time histories
\nextgroupplot[title={standing: postural sway}, timeaxis,
  ylabel={sway [\si{\centi\metre}]}, ymin=-4, ymax=4,
  legend style={at={(0.97,0.03)}, anchor=south east}]
  \addplot[cG!80!black] table[x=t, y expr=100*\thisrow{x}] {results/fig/e_human_standing.dat};
  \addplot[cB] table[x=t, y expr=100*\thisrow{y}] {results/fig/e_human_standing.dat};
  \legend{$x$, $y$}
\nextgroupplot[title={pacing}, timeaxis, ylabel={$x$ [\si{\metre}]}, ymin=-2, ymax=3,
  legend style={at={(0.03,0.97)}, anchor=north west}]
  \addplot[person] table[x=t, y=x] {results/fig/e_human_pacing.dat};
  \addplot[cA] table[x=t, y=rx] {results/fig/e_human_pacing.dat};
  \legend{person, vehicle}
\nextgroupplot[title={sit to stand}, timeaxis, ylabel={height change [\si{\metre}]},
  ymin=-0.55, ymax=0.1, legend style={at={(0.03,0.03)}, anchor=south west}]
  \addplot[person] table[x=t, y=z] {results/fig/e_human_sit_stand.dat};
  \addplot[cA, dashed] table[x=t, y=rzrel] {results/fig/e_human_sit_stand.dat};
  \legend{person, vehicle}
% row 2: top views and the acceleration each motion demands
\nextgroupplot[title={turning around}, topview]
  \addplot[person] table[x=x, y=y] {results/fig/e_human_turnaround.dat};
  \addplot[cA] table[x=rx, y=ry] {results/fig/e_human_turnaround.dat};
  \addplot[standoff] table[x=x, y=y] {results/fig/e_human_turnaround.dat};
  \addplot[stations] table[x=x, y=y] {results/fig/e_human_turnaround.dat};
\nextgroupplot[title={walking a loop, jogging}, topview]
  \addplot[person] table[x=x, y=y] {results/fig/e_human_walk_loop.dat};
  \addplot[cA] table[x=rx, y=ry] {results/fig/e_human_walk_loop.dat};
  \addplot[standoff] table[x=x, y=y] {results/fig/e_human_walk_loop.dat};
  \addplot[stations] table[x=x, y=y] {results/fig/e_human_walk_loop.dat};
\nextgroupplot[title={acceleration demanded}, timeaxis, ymode=log,
  ymin=0.003, ymax=100, ytick={0.01,0.1,1,10,100},
  ylabel={$\|\ddt{\vect r}_{\mathrm{raw}}\|$ [\si{\metre\per\second\squared}]}]
  \addplot[cB] table[x=t, y=acc] {results/fig/e_human_turnaround.dat};
  \addplot[cA] table[x=t, y=acc] {results/fig/e_human_jogging.dat};
  \addplot[cC] table[x=t, y=acc] {results/fig/e_human_pacing.dat};
  \addplot[cD, dashed, domain=0:16, samples=2] {2.5};
  \node[lbl, text=cB, anchor=south] at (axis cs:2.4,22) {turning};
  \node[lbl, text=cA, anchor=south west] at (axis cs:6.4,5.5) {jogging};
  \node[lbl, text=cC!70!black, anchor=north] at (axis cs:9.4,0.35) {pacing};
\end{groupplot}
\end{tikzpicture}

\caption{The six motions flown in \cref{tab:flight}, over one window. Grey is
the person, blue the ideal vehicle position $\vect r_{\mathrm{raw}}$ that holds
the display at the standoff in front of their eyes, and dashed arrows join the
two at intervals. Pacing is mirrored by the vehicle at a fixed offset; sit to
stand is followed in height; on the loop the vehicle runs on an outer circle.
Turning around is different in kind: at each reversal the display must swing
through a half circle of radius $d$ about the person (\cref{prop:turn}). The
last panel shows the price. The raw reference asks for an acceleration an
order of magnitude above the governor's cap $a_{\max}$ (dashed) at every reversal,
jogging's vertical bob alone exceeds it at the stride frequency, and ordinary
walking stays below it.}
\label{fig:missions}
\end{figure}

\begin{table}[t]
\centering
\caption{Missions flown on the hover-sized battery. $\varkappa$ is the mission
overhead of \cref{prop:contraction}; ``session'' is how long the
\SI{\rBattery}{\gram} pack lasts at that power to the \SI{12}{\percent}
reserve. ``Track'' is the peak tracking error over the whole window, start-up
included, to be compared with $\varepsilon=\SI{\rTrackMargin}{\milli\metre}$;
lag is the r.m.s.\ distance from the ideal display position,
\cref{eq:twoerrors}.}
\label{tab:flight}
\footnotesize\setlength{\tabcolsep}{4pt}
\begin{tabular}{@{}lrrrrrrrrr@{}}
\toprule
Mission & Power & $\varkappa$ & Session & Track & Lag rms & Screen tilt & Sat.\ & Envelope & Rotor\\
\midrule
\inputrows{results/tab_flight_rows}
\end{tabular}
\end{table}

The overhead $\varkappa$ ranges from \rKappaMin\ standing to \rKappaMax\ on
the turnaround. The peak tracking error stays below $\varepsilon$ on every
mission (worst \SI{\rTrackWorst}{\milli\metre}), so by
\cref{prop:governor}(iv) the envelope guarantee holds on these runs; the
nearest any rotor disk comes to the eye is \SI{\rClearWorst}{\metre}, on the
turnaround. Disturbance rejection is not what limits the machine indoors: a
\SI{3}{\metre\per\second} breeze with \SI{1.5}{\metre\per\second} gusts, far
stronger than any room, produces \SI{\rBreezeLag}{\milli\metre} of lag and
\ang{\rBreezeTilt} of screen tilt while the airframe leans
\ang{\rBreezeLean}, which is the gimbal doing what it exists for.
\Cref{fig:turn} shows the hardest mission, the turnaround, as time histories.

\begin{figure}[htbp]
\centering
\providecommand{\eDsafeCross}{1.06}
\providecommand{\eDtext}{0.43}
\providecommand{\eDfull}{0.65}
\providecommand{\eLogicalAtCross}{784}
\providecommand{\eArmL}{0.3827}
\providecommand{\eRotorR}{0.2460}
\providecommand{\eSpan}{1.257}
\providecommand{\eReach}{0.629}
\providecommand{\eScreenW}{0.345}
\providecommand{\eScreenH}{0.194}
\providecommand{\eRcp}{0.130}
\providecommand{\eEye}{1.500}
\providecommand{\eZrotor}{1.370}
\providecommand{\eZrotorBelow}{1.630}
\providecommand{\eStandoff}{1.20}
\providecommand{\eDsafe}{0.45}
\providecommand{\eDeltaZ}{0.130}
\providecommand{\eDisplayOrder}{13.3 in laptop panel,15.6 in laptop panel,iPad Pro 13 panel,iPad Pro 13 whole,15.6 in lid assembly,21.5 in monitor panel,24 in monitor panel,13.3 in ultrabook,15.6 in laptop whole}
\providecommand{\eOptSpl}{58}
\providecommand{\eOptMass}{6.0}
\providecommand{\eTurnClear}{0.61}
\providecommand{\eTurnTrack}{64}
\providecommand{\eTrackMargin}{100}
\providecommand{\eTurnPmean}{111.9}
\providecommand{\eHoverBat}{288}
\providecommand{\eReservePct}{12}
\providecommand{\eTethPsrc}{109.6}
\providecommand{\eTethArea}{0.228}
\providecommand{\eTethV}{48}
\providecommand{\eTethL}{5}
\providecommand{\eTethDrop}{5}
%
\begin{tikzpicture}
\begin{groupplot}[
  group style={group size=1 by 4, vertical sep=0.22cm,
               xlabels at=edge bottom, xticklabels at=edge bottom},
  paperplot, height=3.2cm, xmin=0, xmax=16,
  xlabel={time $t$ [\si{\second}]},
  ylabel style={align=center, font=\footnotesize}]
\nextgroupplot[ylabel={$e_{\mathrm{track}}$\\ {[}\si{\milli\metre}{]}}, ymin=0, ymax=125]
  \fill[cG!12] (axis cs:0,0) rectangle (axis cs:3,125);
  \addplot[cA] table[x=t, y=track] {results/fig/e_turn.dat};
  \addplot[cD, dashed, domain=0:16, samples=2] {\eTrackMargin};
  \node[lbl, text=cD, anchor=south east] at (axis cs:15.9,\eTrackMargin) {$\varepsilon$};
  \node[lbl, text=cG!80!black, anchor=north west] at (axis cs:0.1,90) {settling};
\nextgroupplot[ylabel={$e_{\mathrm{lag}}$\\ {[}\si{\metre}{]}}, ymin=0]
  \fill[cG!12] (axis cs:0,0) rectangle (axis cs:3,4);
  \addplot[cB] table[x=t, y expr=\thisrow{lag}/1000] {results/fig/e_turn.dat};
\nextgroupplot[ylabel={rotor to eye\\ {[}\si{\metre}{]}}, ymin=0.3, ymax=1.5]
  \fill[cG!12] (axis cs:0,0.3) rectangle (axis cs:3,1.5);
  \addplot[cC!70!black] table[x=t, y=clear] {results/fig/e_turn.dat};
  \addplot[cD, dashed, domain=0:16, samples=2] {\eDsafe};
  \node[lbl, text=cD, anchor=south] at (axis cs:4.6,\eDsafe) {$d_{\mathrm{safe}}$};
\nextgroupplot[ylabel={power\\ {[}\si{\watt}{]}}]
  \fill[cG!12] (axis cs:0,90) rectangle (axis cs:3,140);
  \addplot[cE] table[x=t, y=P] {results/fig/e_turn.dat};
\end{groupplot}
\end{tikzpicture}

\caption{The turnaround flown with the governor. (a) The tracking error
against the governed reference stays below the margin $\varepsilon$, which is
the hypothesis of \cref{prop:governor}(iv) checked on this run. (b) The lag
behind the ideal display position grows to about \SI{2}{\metre} at each
reversal: the governor lets the display fall behind rather than swing it round
the head at the acceleration of \cref{fig:missions}. (c) The distance from the
eye to the nearest rotor disk never approaches $d_{\mathrm{safe}}$. (d) The
electrical power, whose window mean after the settling interval (grey) is the
$\bar P$ of the energy integral; the brief spikes are the rotor speeds
responding to the jerk of the reference at the start and end of each swing.}
\label{fig:turn}
\end{figure}

What the table also shows is that a battery sized for hover is not sized for
the mission: on every moving mission the pack lands below its reserve, at
\SI{7}{\percent} on the turnaround. That is the reason the coupled loop exists.

\FloatBarrier
\subsection{The fully coupled loop, following the user}

The design of \cref{tab:design} was sized for hover and then flown. That is
not \cref{prob:codesign}. The coupled problem sizes the battery around the
motion actually performed: guess $\mbat$, close frame and propulsion around it
by \cref{prop:inner}, rebuild the vehicle including its inertia tensor and
rotor constants, fly the mission with the governor and keep-out, integrate the
electrical power, and ask what battery that mission required, by
\cref{eq:batreq}. \Cref{tab:coupled} gives the result for each motion. Every
row converged in \rCoupItLo\ battery updates by the secant iteration
\cref{eq:secant} followed by two verification flights, \rCoupSimLo\ full
6-DOF simulations in all, and every row's final flight met the reserve, the
peak-power requirement, the rotor clearance and the readability limit.

\begin{table}[t]
\centering
\caption{The fully coupled closure \cref{eq:residual} on the final design, one
row per motion. ``Hover sizing'' is the first-principles closure of
\cref{eq:fpclosure} with the mission assumed to be hover, i.e.\ the design of
\cref{tab:design}. ``Updates/sims'' counts battery updates and total
simulations including the verification flights; ``reserve'' is the state of
charge left at $t_f$ on the final flight; ``rotor'' the minimum eye-to-disk
distance on it; ``feasible'' whether that flight met every path constraint.}
\label{tab:coupled}
\footnotesize\setlength{\tabcolsep}{4pt}
\begin{tabular}{@{}lrrrrrrrrl@{}}
\toprule
Motion & Gross & Battery & vs.\ hover & Mean power & $\varkappa$ & Upd./sims & Reserve & Rotor & Feasible\\
\midrule
hover sizing & \SI{\rGross}{\kilo\gram} & \SI{\rBattery}{\gram} & n/a & \SI{\rPower}{\watt} & n/a & n/a & \SI{12.0}{\percent} & n/a & n/a\\
\midrule
\inputrows{results/tab_coupled_rows}
\end{tabular}
\end{table}

\begin{figure}[tbp]
\centering
\providecommand{\eDsafeCross}{1.06}
\providecommand{\eDtext}{0.43}
\providecommand{\eDfull}{0.65}
\providecommand{\eLogicalAtCross}{784}
\providecommand{\eArmL}{0.3827}
\providecommand{\eRotorR}{0.2460}
\providecommand{\eSpan}{1.257}
\providecommand{\eReach}{0.629}
\providecommand{\eScreenW}{0.345}
\providecommand{\eScreenH}{0.194}
\providecommand{\eRcp}{0.130}
\providecommand{\eEye}{1.500}
\providecommand{\eZrotor}{1.370}
\providecommand{\eZrotorBelow}{1.630}
\providecommand{\eStandoff}{1.20}
\providecommand{\eDsafe}{0.45}
\providecommand{\eDeltaZ}{0.130}
\providecommand{\eDisplayOrder}{13.3 in laptop panel,15.6 in laptop panel,iPad Pro 13 panel,iPad Pro 13 whole,15.6 in lid assembly,21.5 in monitor panel,24 in monitor panel,13.3 in ultrabook,15.6 in laptop whole}
\providecommand{\eOptSpl}{58}
\providecommand{\eOptMass}{6.0}
\providecommand{\eTurnClear}{0.61}
\providecommand{\eTurnTrack}{64}
\providecommand{\eTrackMargin}{100}
\providecommand{\eTurnPmean}{111.9}
\providecommand{\eHoverBat}{288}
\providecommand{\eReservePct}{12}
\providecommand{\eTethPsrc}{109.6}
\providecommand{\eTethArea}{0.228}
\providecommand{\eTethV}{48}
\providecommand{\eTethL}{5}
\providecommand{\eTethDrop}{5}
%
\begin{tikzpicture}
\begin{groupplot}[
  group style={group size=2 by 1, horizontal sep=1.6cm},
  halfplot, height=5.2cm, ybar, /pgf/bar width=13pt,
  symbolic x coords={standing,pacing,loop,jogging,turning},
  xtick=data, enlarge x limits=0.14,
  x tick label style={font=\footnotesize, rotate=25, anchor=north east}]
\nextgroupplot[ylabel={battery the mission needs [\si{\gram}]}, ymin=270, ymax=320,
  title={closed around the flown mission}]
  \addplot[fill=cA!70, draw=cA] table[x=mission, y=mbat] {results/fig/e_coupled.dat};
  \draw[cG, dashed, line width=0.8pt]
    ({rel axis cs:0,0} |- {axis cs:standing,\eHoverBat}) --
    ({rel axis cs:1,0} |- {axis cs:standing,\eHoverBat});
  \node[lbl, text=cG!80!black, anchor=north west] at (rel axis cs:0.02,0.97)
    {dashed: hover sizing, \eHoverBat\ g};
\nextgroupplot[ylabel={charge left with \eHoverBat\ g [\si{\percent}]}, ymin=0, ymax=14,
  title={hover-sized battery, flown}]
  \addplot[fill=cB!65, draw=cB] table[x=mission, y=soc] {results/fig/e_coupled.dat};
  \draw[cD, dashed, line width=0.8pt]
    ({rel axis cs:0,0} |- {axis cs:standing,\eReservePct}) --
    ({rel axis cs:1,0} |- {axis cs:standing,\eReservePct})
    node[pos=0.98, below left, lbl, text=cD] {reserve \eReservePct\,\%};
\end{groupplot}
\end{tikzpicture}

\caption{The coupled closure against the hover sizing. Left: the battery each
motion needs when \cref{eq:residual} is solved around the flown mission; the
dashed line is the \SI{\eHoverBat}{\gram} of the hover sizing. Right: the
charge left at $t_f$ when that hover-sized pack is flown through each motion
instead. Every moving mission lands below the \SI{\eReservePct}{\percent}
reserve. The two panels are the same fact from either side: the right is the
violation, the left is its cure.}
\label{fig:coupled}
\end{figure}

\Cref{fig:coupled} shows the table's battery column against what the hover
sizing would have done. Three observations follow.

\emph{The standing row reproduces the hover sizing} to within
\SI{\rCoupStandBat}{\gram} of battery. This is a consistency check, not a
validation: the simulation and the closure share the rotor model of
\cref{sec:blade}, and agreement between them says that the substitution of
\cref{mod:battery} at $\varkappa=1$ is implemented correctly, nothing more.

\emph{Following a person is cheap; turning with them is not.} Ordinary motion
costs between \SI{\rCoupOrdLo}{\gram} and \SI{\rCoupOrdHi}{\gram} of battery.
Repeated turnarounds cost \SI{\rCoupTurnBat}{\gram} of battery and
\SI{\rCoupTurnGross}{\gram} of gross mass, almost all of it the mission
overhead $\varkappa=\rCoupTurnKappa$ propagated through the fixed point: by
\cref{eq:kappaclosure} the overhead multiplies $\beta$, and here it moves the
design \SI{\rCoupTurnUp}{\percent} further up the closure curve. Sizing the
battery for the most demanding motion is a worst-case operating policy, and
it is the one adopted here. For a defined session with known phases the
energy is the time-weighted sum over the phases, with the peak-power
requirement enforced separately by \cref{eq:batreq}; a session of occasional
turns needs less than \cref{tab:coupled}'s last row, and one of continuous
turning needs exactly it.

\emph{The inertia coupling is weak for this layout.} The cycle \cref{eq:cycle2}
runs through $J$, but the battery sits near the centre of mass, so $J_{yy}$
moves only from \num{\rCoupJlo} to \SI{\rCoupJhi}{\kilo\gram\metre\squared}
between the lightest and heaviest converged designs. The coupling that matters
is the energy one; the dynamic one is present and small.

\paragraph{Dynamic clearance.} The static clearance of \cref{tab:design},
\SI{\rTipGap}{\metre}, is not the clearance flown. During the turnaround the
governed reference is held outside $\varrho = d_{\mathrm{safe}}+\mathcal{R}+
\varepsilon$ of the head by \cref{eq:barrier}, the vehicle tracks it to within
\SI{\rTurnTrack}{\milli\metre}, and the nearest rotor disk passes the eye at
\SI{\rTurnClear}{\metre} (envelope gap \SI{\rTurnEnv}{\metre}). An earlier
version of the governor, with a keep-out on position only and a velocity test
that ignored the head's own motion, let the vehicle inside the safe distance
by about \SI{4}{\centi\metre} on this mission; the margin $\varepsilon$ and
the moving-boundary barrier are what closed that gap, and
\cref{prop:governor}(iv) states exactly the hypothesis under which they do.

\FloatBarrier
\subsection{The tether}

Taking the battery off the aircraft removes the mission time from the closure.
What replaces it is a conductor whose mass follows from Ohm's law: at bus
voltage $V$ delivering power $P$ over length $L$ with fractional drop
$\delta$,
\begin{equation}
  I = \frac{P}{V},
  \quad
  R_{\mathrm{cond}} = \frac{2\rho_e L}{A_c},
  \quad
  IR_{\mathrm{cond}}\le\delta V
  \;\Longrightarrow\;
  A_c \ge \frac{2\rho_e L P}{\delta V^2},
  \quad
  m_c = 2\rho_m L A_c ,
  \label{eq:tether}
\end{equation}
with $A_c$ taken as the largest of the resistive requirement, the continuous
current density limit $I/J_{\max}$ with $J_{\max}=\SI{\rTethJmax}{\ampere\per\milli\metre\squared}$,
and a handling floor of \SI{\rTethAreaMin}{\milli\metre\squared}.
\Cref{fig:tether} shows how the three requirements trade over power.

\begin{figure}[tbp]
\centering
\providecommand{\eDsafeCross}{1.06}
\providecommand{\eDtext}{0.43}
\providecommand{\eDfull}{0.65}
\providecommand{\eLogicalAtCross}{784}
\providecommand{\eArmL}{0.3827}
\providecommand{\eRotorR}{0.2460}
\providecommand{\eSpan}{1.257}
\providecommand{\eReach}{0.629}
\providecommand{\eScreenW}{0.345}
\providecommand{\eScreenH}{0.194}
\providecommand{\eRcp}{0.130}
\providecommand{\eEye}{1.500}
\providecommand{\eZrotor}{1.370}
\providecommand{\eZrotorBelow}{1.630}
\providecommand{\eStandoff}{1.20}
\providecommand{\eDsafe}{0.45}
\providecommand{\eDeltaZ}{0.130}
\providecommand{\eDisplayOrder}{13.3 in laptop panel,15.6 in laptop panel,iPad Pro 13 panel,iPad Pro 13 whole,15.6 in lid assembly,21.5 in monitor panel,24 in monitor panel,13.3 in ultrabook,15.6 in laptop whole}
\providecommand{\eOptSpl}{58}
\providecommand{\eOptMass}{6.0}
\providecommand{\eTurnClear}{0.61}
\providecommand{\eTurnTrack}{64}
\providecommand{\eTrackMargin}{100}
\providecommand{\eTurnPmean}{111.9}
\providecommand{\eHoverBat}{288}
\providecommand{\eReservePct}{12}
\providecommand{\eTethPsrc}{109.6}
\providecommand{\eTethArea}{0.228}
\providecommand{\eTethV}{48}
\providecommand{\eTethL}{5}
\providecommand{\eTethDrop}{5}
%
\begin{tikzpicture}
\begin{axis}[paperplot, height=5.4cm, xmin=0, xmax=400, ymin=0, ymax=0.9,
  xlabel={power drawn at the source $P$ [\si{\watt}]},
  ylabel={conductor area per leg [\si{\milli\metre\squared}]},
  legend style={at={(0.02,0.98)}, anchor=north west}]
  \addplot[cA] table[x=P, y=drop] {results/fig/e_tether.dat};
  \addplot[cB] table[x=P, y=amp] {results/fig/e_tether.dat};
  \addplot[cG, dashed] table[x=P, y=hand] {results/fig/e_tether.dat};
  \addplot[black, line width=2.2pt, opacity=0.25] table[x=P, y=req] {results/fig/e_tether.dat};
  \addplot[only marks, mark=*, mark size=2.2pt, cD]
    coordinates {(\eTethPsrc,\eTethArea)};
  \draw[cD, line width=0.5pt] (axis cs:\eTethPsrc,\eTethArea) -- (axis cs:190,0.05);
  \node[lbl, text=cD, anchor=west, fill=white, inner sep=1.5pt] at (axis cs:190,0.05)
    {design: \SI{\eTethArea}{\milli\metre\squared} at \SI{\eTethPsrc}{\watt}};
  \legend{voltage drop $\le\SI{\eTethDrop}{\percent}$,
          current density $\le J_{\max}$,
          handling floor,
          required: the largest}
\end{axis}
\end{tikzpicture}

\caption{Conductor area per leg for a \SI{\eTethL}{\metre} lead at
\SI{\eTethV}{\volt}, against the power drawn at the source. Each requirement of
\cref{eq:tether} is linear in $P$ or constant, and the conductor is sized by
their maximum (grey band): the handling floor at low power, the current
density limit above about \SI{60}{\watt}. At this length and voltage the
voltage-drop requirement, the only one that grows with $L^2$, never binds.
The design point lies on the ampacity line, so the carried conductor mass is
proportional to power, which is the $m^{3/2}$ term that keeps the fold.}
\label{fig:tether}
\end{figure}

\begin{observation}
\Cref{eq:tether} contains no mission time, so endurance ceases to be a design
variable and becomes an operating choice. It does not delete the fold. A
carried conductor sized by resistance has mass proportional to $P\propto
m^{3/2}$, which is the same term as the battery's with a coefficient smaller by
orders of magnitude at these lengths and voltages: the fold moves far away, it
does not vanish. The $m^{3/2}$ term is absent only when the carried cable mass
stops depending on power, because the tether is supported from above or
because a gauge floor binds. For the case below the lumped closure is
gauge-limited (\rTethGaugeLimited), which is why its $\beta_{\mathrm{tether}}$
is zero.
\end{observation}

Redesigning the aircraft for a \SI{\rTethLen}{\metre} lead at
\SI{\rTethV}{\volt}, with the same component models as the battery closure,
gives \SI{\rTethGross}{\kilo\gram} at \SI{\rTethPower}{\watt} and
\SI{\rTethEar}{\decibel A} at the ear. This is not the battery design with the
battery removed: propulsion and structure are resized around the lighter
aircraft, and it still carries a \SI{\rTethReserve}{\gram} landing reserve,
sized by the peak power it must deliver (\rTethResLimit-limited) rather than by
energy. The current is \SI{\rTethCurrent}{\ampere}; the conductor is set by
\rTethSizedBy\ at \SI{\rTethArea}{\milli\metre\squared} per leg and
\SI{\rTethMass}{\gram} for the lead. \SI{\rTethV}{\volt} remains below the
\SI{60}{\volt} DC safety-extra-low-voltage threshold, and the design
\rTethFeasible\ the convergence and current-density checks the engine applies.
